\section{Introduction}
\label{sec:intro}

An online log processor may need to identify an event and choose a follow-up
check before a response deadline. A correct answer arriving after that
deadline can still help later analysis, but cannot support the intended
timely decision. This makes response latency a different objective from
processing as many requests per second as possible. The operational goal
is to maximize the fraction of problems resolved correctly before their
deadline, while retaining the required quality standard.

Autoregressive (AR) decoding generates tokens, the model's units of text,
one after another. Block diffusion instead refines a fixed-size block of
token positions, called a \emph{canvas}, through several passes before
moving to the next block. It can therefore produce a structured response
with fewer sequential decoding steps. Figure~\ref{fig:shapes} connects this
intuition to the measured latency gap and conditional output-length model.

Our primary comparison is Qwen3-30B-A3B (AR-30B)~\cite{qwen3} versus
DiffusionGemma-26B-A4B (DG)~\cite{diffusiongemma}. These similarly sized mixture-of-experts
models activate only a subset of their parameters for each token:
AR has approximately 30 billion total and 3 billion active parameters,
while DG has 26 billion total and 4 billion active parameters.
They are comparable in scale, not identical in training or capability.
Experiments with dense Qwen2.5-3B (AR-3B) and the separate LLaDA-8B
diffusion model provide supporting evidence. Section~\ref{sec:survives}
uses one explicitly labeled AR-3B timing illustration; details of these
secondary experiments remain in Appendix~\ref{app:secondary}.

\noindent\begin{minipage}{\linewidth}
\paragraph{Three takeaways.}
\begin{enumerate}
\item \textbf{Diffusion returns sooner in the measured serial workload.}
The primary pair has median generation times of approximately 924\,ms
for AR and 380\,ms for DG, a $2.4\times$ ratio.
\item \textbf{The latency gap changes which responses fit a time budget.}
Within 500\,ms, 99.8\% of DG generations finish in the saved runs,
compared with none from this AR configuration.
\item \textbf{Quality and concurrency determine whether that speed is useful.}
On 300 shared events, nine-field F1 is 0.870 for AR and 0.852 for DG.
The scores are close, with a small DG deficit; AR also gains
more throughput from batching.
\end{enumerate}
\end{minipage}

\paragraph{Task and measures.}
Each BlueGene/L (BGL) system log line from Loghub~\cite{loghub} is converted to a JavaScript Object
Notation (JSON) record with twelve named fields. We evaluate nine fields:
seven structural fields and two anomaly fields. The other three describe
events, affected entities, and actions, whose semantic quality is not
independently labeled. Field-F1 measures token overlap with parser-derived
reference values; it is the harmonic mean of precision and recall, averaged
equally over scored fields. Schema validity checks required fields,
types, and allowed values. Neither guarantees a correct downstream decision.

We use saved serial runs on an NVIDIA A100 graphics processing unit (GPU)
with 80\,GB memory. Both models use vLLM serving software and greedy
decoding, without random sampling. DG uses eight refinement passes
($K=8$) and a 256-position canvas ($C=256$); AR has a 256-token maximum
output budget. The latency comparison pools three AR runs on
300 events and one DG run on 2,000 events. Paired quality uses the first
300 shared events, not the pooled timing samples.
Preparation, scoring, and conflicting DG software-version metadata
are documented in Appendix~\ref{app:protocol}.

\paragraph{Roadmap.}
We explain the cost models, compare the primary pair's quality and
generation-time budgets, then discuss batching and how latency savings
can accumulate across dependent actions. The appendices retain protocol details, smaller-model
studies, and exploratory results without mixing them into the main comparison.
